\section{领域三：Claude Code 配置与工作流（20\%）}

这个领域考察的是你是否不仅会\textbf{使用} Claude Code，还能为团队\textbf{配置}它。核心概念包括 CLAUDE.md 层级体系、规则目录、plan mode 与 direct execution 的选择。

\subsection{CLAUDE.md 三级层次结构}

\begin{importantbox}{CLAUDE.md 的三个层级}
\begin{enumerate}[leftmargin=1.5em]
    \item \textbf{User-level}：\texttt{\textasciitilde/.claude/CLAUDE.md}——个人配置，不受版本控制，不会与团队共享
    \item \textbf{Project-level}：\texttt{.claude/CLAUDE.md}（或项目根目录的 \texttt{CLAUDE.md}）——项目级配置，随代码库版本控制
    \item \textbf{Directory-level}：子目录中的 \texttt{CLAUDE.md}——仅对该目录及其子目录生效
\end{enumerate}
\end{importantbox}

\begin{warningbox}{考试常见陷阱：User-level 配置不共享}
如果团队成员缺少某些指令，很可能是因为这些指令被放在了 user-level 的 \texttt{CLAUDE.md} 中。由于该文件不受版本控制，其他团队成员不会看到这些配置。团队共享的配置必须放在 project-level。
\end{warningbox}

\subsection{路径特定规则（Path-Specific Rules）}

\begin{knowledgebox}{.claude/rules/ 目录的 Glob 模式}
\texttt{.claude/rules/} 目录支持 YAML frontmatter 中的 glob 模式，例如 \texttt{**/*.test.tsx} 可以将特定约定应用于整个代码库中匹配的文件。这是 directory-level 的 \texttt{CLAUDE.md} 做不到的，因为后者只对所在目录及子目录生效，无法跨目录匹配文件路径模式。
\end{knowledgebox}

\subsection{Plan Mode vs Direct Execution}

\begin{table}[H]
\centering
\begin{tabular}{lp{9cm}}
\toprule
\textbf{模式} & \textbf{适用场景} \\
\midrule
Plan Mode & 单体应用重构、多文件迁移、架构决策等需要全局视角的任务 \\
Direct Execution & 单文件 bug 修复、单个验证检查等范围明确的小任务 \\
\bottomrule
\end{tabular}
\caption{Plan Mode 与 Direct Execution 的适用场景}
\end{table}

\subsection{其他关键概念}

需要掌握的其他 Claude Code 概念：
\begin{itemize}[leftmargin=1.5em]
    \item \textbf{\texttt{context: fork}}：在 skill frontmatter 中使用，隔离冗长输出，避免污染主会话上下文
    \item \textbf{\texttt{-p} flag}：非交互式模式，用于 CI/CD 管道
    \item \textbf{独立审查实例}：用一个新的 Claude Code 实例来审查代码，比在同一个 session 中自我审查效果更好
\end{itemize}

\subsection{推荐实践项目}

\begin{knowledgebox}{动手练习：完整 Claude Code 配置}
创建一个项目，包含：
\begin{itemize}[leftmargin=1.5em]
    \item 完整的 CLAUDE.md 三级层次结构
    \item \texttt{.claude/rules/} 中带 glob 模式的规则文件
    \item 一个使用 \texttt{context: fork} 的 skill
    \item \texttt{.mcp.json} 中配置带环境变量展开的 MCP server
\end{itemize}
然后分别用 plan mode 处理多文件重构，用 direct execution 处理单个 bug 修复，体会两者的区别。
\end{knowledgebox}

\subsection{本章小结}

Claude Code 的团队配置能力是架构师的分水岭。核心要点：理解 CLAUDE.md 的三级层次并将共享配置放在 project-level；利用 \texttt{.claude/rules/} 的 glob 模式实现跨目录的路径特定规则；根据任务范围选择 plan mode 或 direct execution。

